\documentclass[11pt,a4paper]{article}
% lesson-preamble.tex -- Shared preamble for all Phase 00 lessons
% Usage: % lesson-preamble.tex -- Shared preamble for all Phase 00 lessons
% Usage: % lesson-preamble.tex -- Shared preamble for all Phase 00 lessons
% Usage: \input{../lesson-preamble} (from subdomain directories)

% ---- Encoding and fonts ----
\usepackage[utf8]{inputenc}
\usepackage[T1]{fontenc}

% ---- Mathematics ----
\usepackage{amsmath,amssymb,amsthm}
\usepackage{mathtools}

% ---- Page geometry ----
\usepackage[margin=1in,headheight=14pt]{geometry}

% ---- Hyperlinks ----
\usepackage[colorlinks=true,linkcolor=red,citecolor=blue,urlcolor=blue]{hyperref}

% ---- Lists ----
\usepackage{enumitem}

% ---- Colors and boxes ----
\usepackage{xcolor}
\usepackage[theorems,skins,breakable]{tcolorbox}

% ---- Header/Footer ----
\usepackage{fancyhdr}
\renewcommand{\headrulewidth}{0.4pt}

% ---- Tables ----
\usepackage{booktabs}

% ---- Bibliography ----
\usepackage{natbib}
\bibliographystyle{plainnat}

% --------------------------------------------------------------------------
% Theorem environments
% --------------------------------------------------------------------------
\theoremstyle{plain}
\newtheorem{theorem}{Theorem}[section]
\newtheorem{lemma}[theorem]{Lemma}
\newtheorem{proposition}[theorem]{Proposition}
\newtheorem{corollary}[theorem]{Corollary}

\theoremstyle{definition}
\newtheorem{definition}[theorem]{Definition}
\newtheorem{example}[theorem]{Example}
\newtheorem{exercise}[theorem]{Exercise}
\newtheorem{assumption}[theorem]{Assumption}

\theoremstyle{remark}
\newtheorem{remark}[theorem]{Remark}

% --------------------------------------------------------------------------
% Colored boxes
% --------------------------------------------------------------------------
\newtcolorbox{keyresult}[1][]{%
  colback=green!5!white,
  colframe=green!50!black,
  fonttitle=\bfseries,
  title={Key Result},
  breakable,
  #1
}

\newtcolorbox{rlconnection}[1][]{%
  colback=blue!5!white,
  colframe=blue!50!black,
  fonttitle=\bfseries,
  title={RL/ML Connection},
  breakable,
  #1
}

\newtcolorbox{warning}[1][]{%
  colback=red!5!white,
  colframe=red!50!black,
  fonttitle=\bfseries,
  title={Warning},
  breakable,
  #1
}

% --------------------------------------------------------------------------
% Custom commands -- number sets
% --------------------------------------------------------------------------
\newcommand{\R}{\mathbb{R}}
\newcommand{\N}{\mathbb{N}}
\newcommand{\Q}{\mathbb{Q}}
\newcommand{\Z}{\mathbb{Z}}
\newcommand{\C}{\mathbb{C}}
\newcommand{\F}{\mathbb{F}}

% --------------------------------------------------------------------------
% Custom commands -- probability and expectation
% --------------------------------------------------------------------------
\newcommand{\E}{\mathbb{E}}
\newcommand{\Prob}{\mathbb{P}}
\newcommand{\Var}{\operatorname{Var}}
\newcommand{\Cov}{\operatorname{Cov}}
\newcommand{\indicator}{\mathbf{1}}

% --------------------------------------------------------------------------
% Custom commands -- convergence arrows
% --------------------------------------------------------------------------
\newcommand{\cas}{\xrightarrow{\text{a.s.}}}
\newcommand{\cprob}{\xrightarrow{\Prob}}
\newcommand{\clp}[1]{\xrightarrow{L^{#1}}}
\newcommand{\cdist}{\xrightarrow{d}}

% --------------------------------------------------------------------------
% Custom commands -- common operators
% --------------------------------------------------------------------------
\DeclareMathOperator{\dom}{dom}
\DeclareMathOperator{\epi}{epi}
\DeclareMathOperator{\conv}{conv}
\DeclareMathOperator{\relint}{relint}
\DeclareMathOperator{\aff}{aff}
\DeclareMathOperator{\cl}{cl}
\DeclareMathOperator{\interior}{int}
\DeclareMathOperator{\tr}{tr}
\DeclareMathOperator{\diag}{diag}
\DeclareMathOperator{\rank}{rank}
\DeclareMathOperator{\sgn}{sgn}
\DeclareMathOperator{\supp}{supp}
\DeclareMathOperator{\argmin}{arg\,min}
\DeclareMathOperator{\argmax}{arg\,max}
\DeclareMathOperator{\prox}{prox}
\DeclareMathOperator{\dist}{dist}
\DeclareMathOperator{\diam}{diam}
\DeclareMathOperator{\Span}{span}

% --------------------------------------------------------------------------
% Custom commands -- linear algebra operators
% --------------------------------------------------------------------------
\DeclareMathOperator{\nullity}{nullity}
\DeclareMathOperator{\range}{range}
\DeclareMathOperator{\nullspace}{null}
\DeclareMathOperator{\proj}{proj}

% --------------------------------------------------------------------------
% Custom commands -- measure theory
% --------------------------------------------------------------------------
\newcommand{\powerset}{\mathcal{P}}
\newcommand{\Borel}{\mathcal{B}}
 (from subdomain directories)

% ---- Encoding and fonts ----
\usepackage[utf8]{inputenc}
\usepackage[T1]{fontenc}

% ---- Mathematics ----
\usepackage{amsmath,amssymb,amsthm}
\usepackage{mathtools}

% ---- Page geometry ----
\usepackage[margin=1in,headheight=14pt]{geometry}

% ---- Hyperlinks ----
\usepackage[colorlinks=true,linkcolor=red,citecolor=blue,urlcolor=blue]{hyperref}

% ---- Lists ----
\usepackage{enumitem}

% ---- Colors and boxes ----
\usepackage{xcolor}
\usepackage[theorems,skins,breakable]{tcolorbox}

% ---- Header/Footer ----
\usepackage{fancyhdr}
\renewcommand{\headrulewidth}{0.4pt}

% ---- Tables ----
\usepackage{booktabs}

% ---- Bibliography ----
\usepackage{natbib}
\bibliographystyle{plainnat}

% --------------------------------------------------------------------------
% Theorem environments
% --------------------------------------------------------------------------
\theoremstyle{plain}
\newtheorem{theorem}{Theorem}[section]
\newtheorem{lemma}[theorem]{Lemma}
\newtheorem{proposition}[theorem]{Proposition}
\newtheorem{corollary}[theorem]{Corollary}

\theoremstyle{definition}
\newtheorem{definition}[theorem]{Definition}
\newtheorem{example}[theorem]{Example}
\newtheorem{exercise}[theorem]{Exercise}
\newtheorem{assumption}[theorem]{Assumption}

\theoremstyle{remark}
\newtheorem{remark}[theorem]{Remark}

% --------------------------------------------------------------------------
% Colored boxes
% --------------------------------------------------------------------------
\newtcolorbox{keyresult}[1][]{%
  colback=green!5!white,
  colframe=green!50!black,
  fonttitle=\bfseries,
  title={Key Result},
  breakable,
  #1
}

\newtcolorbox{rlconnection}[1][]{%
  colback=blue!5!white,
  colframe=blue!50!black,
  fonttitle=\bfseries,
  title={RL/ML Connection},
  breakable,
  #1
}

\newtcolorbox{warning}[1][]{%
  colback=red!5!white,
  colframe=red!50!black,
  fonttitle=\bfseries,
  title={Warning},
  breakable,
  #1
}

% --------------------------------------------------------------------------
% Custom commands -- number sets
% --------------------------------------------------------------------------
\newcommand{\R}{\mathbb{R}}
\newcommand{\N}{\mathbb{N}}
\newcommand{\Q}{\mathbb{Q}}
\newcommand{\Z}{\mathbb{Z}}
\newcommand{\C}{\mathbb{C}}
\newcommand{\F}{\mathbb{F}}

% --------------------------------------------------------------------------
% Custom commands -- probability and expectation
% --------------------------------------------------------------------------
\newcommand{\E}{\mathbb{E}}
\newcommand{\Prob}{\mathbb{P}}
\newcommand{\Var}{\operatorname{Var}}
\newcommand{\Cov}{\operatorname{Cov}}
\newcommand{\indicator}{\mathbf{1}}

% --------------------------------------------------------------------------
% Custom commands -- convergence arrows
% --------------------------------------------------------------------------
\newcommand{\cas}{\xrightarrow{\text{a.s.}}}
\newcommand{\cprob}{\xrightarrow{\Prob}}
\newcommand{\clp}[1]{\xrightarrow{L^{#1}}}
\newcommand{\cdist}{\xrightarrow{d}}

% --------------------------------------------------------------------------
% Custom commands -- common operators
% --------------------------------------------------------------------------
\DeclareMathOperator{\dom}{dom}
\DeclareMathOperator{\epi}{epi}
\DeclareMathOperator{\conv}{conv}
\DeclareMathOperator{\relint}{relint}
\DeclareMathOperator{\aff}{aff}
\DeclareMathOperator{\cl}{cl}
\DeclareMathOperator{\interior}{int}
\DeclareMathOperator{\tr}{tr}
\DeclareMathOperator{\diag}{diag}
\DeclareMathOperator{\rank}{rank}
\DeclareMathOperator{\sgn}{sgn}
\DeclareMathOperator{\supp}{supp}
\DeclareMathOperator{\argmin}{arg\,min}
\DeclareMathOperator{\argmax}{arg\,max}
\DeclareMathOperator{\prox}{prox}
\DeclareMathOperator{\dist}{dist}
\DeclareMathOperator{\diam}{diam}
\DeclareMathOperator{\Span}{span}

% --------------------------------------------------------------------------
% Custom commands -- linear algebra operators
% --------------------------------------------------------------------------
\DeclareMathOperator{\nullity}{nullity}
\DeclareMathOperator{\range}{range}
\DeclareMathOperator{\nullspace}{null}
\DeclareMathOperator{\proj}{proj}

% --------------------------------------------------------------------------
% Custom commands -- measure theory
% --------------------------------------------------------------------------
\newcommand{\powerset}{\mathcal{P}}
\newcommand{\Borel}{\mathcal{B}}
 (from subdomain directories)

% ---- Encoding and fonts ----
\usepackage[utf8]{inputenc}
\usepackage[T1]{fontenc}

% ---- Mathematics ----
\usepackage{amsmath,amssymb,amsthm}
\usepackage{mathtools}

% ---- Page geometry ----
\usepackage[margin=1in,headheight=14pt]{geometry}

% ---- Hyperlinks ----
\usepackage[colorlinks=true,linkcolor=red,citecolor=blue,urlcolor=blue]{hyperref}

% ---- Lists ----
\usepackage{enumitem}

% ---- Colors and boxes ----
\usepackage{xcolor}
\usepackage[theorems,skins,breakable]{tcolorbox}

% ---- Header/Footer ----
\usepackage{fancyhdr}
\renewcommand{\headrulewidth}{0.4pt}

% ---- Tables ----
\usepackage{booktabs}

% ---- Bibliography ----
\usepackage{natbib}
\bibliographystyle{plainnat}

% --------------------------------------------------------------------------
% Theorem environments
% --------------------------------------------------------------------------
\theoremstyle{plain}
\newtheorem{theorem}{Theorem}[section]
\newtheorem{lemma}[theorem]{Lemma}
\newtheorem{proposition}[theorem]{Proposition}
\newtheorem{corollary}[theorem]{Corollary}

\theoremstyle{definition}
\newtheorem{definition}[theorem]{Definition}
\newtheorem{example}[theorem]{Example}
\newtheorem{exercise}[theorem]{Exercise}
\newtheorem{assumption}[theorem]{Assumption}

\theoremstyle{remark}
\newtheorem{remark}[theorem]{Remark}

% --------------------------------------------------------------------------
% Colored boxes
% --------------------------------------------------------------------------
\newtcolorbox{keyresult}[1][]{%
  colback=green!5!white,
  colframe=green!50!black,
  fonttitle=\bfseries,
  title={Key Result},
  breakable,
  #1
}

\newtcolorbox{rlconnection}[1][]{%
  colback=blue!5!white,
  colframe=blue!50!black,
  fonttitle=\bfseries,
  title={RL/ML Connection},
  breakable,
  #1
}

\newtcolorbox{warning}[1][]{%
  colback=red!5!white,
  colframe=red!50!black,
  fonttitle=\bfseries,
  title={Warning},
  breakable,
  #1
}

% --------------------------------------------------------------------------
% Custom commands -- number sets
% --------------------------------------------------------------------------
\newcommand{\R}{\mathbb{R}}
\newcommand{\N}{\mathbb{N}}
\newcommand{\Q}{\mathbb{Q}}
\newcommand{\Z}{\mathbb{Z}}
\newcommand{\C}{\mathbb{C}}
\newcommand{\F}{\mathbb{F}}

% --------------------------------------------------------------------------
% Custom commands -- probability and expectation
% --------------------------------------------------------------------------
\newcommand{\E}{\mathbb{E}}
\newcommand{\Prob}{\mathbb{P}}
\newcommand{\Var}{\operatorname{Var}}
\newcommand{\Cov}{\operatorname{Cov}}
\newcommand{\indicator}{\mathbf{1}}

% --------------------------------------------------------------------------
% Custom commands -- convergence arrows
% --------------------------------------------------------------------------
\newcommand{\cas}{\xrightarrow{\text{a.s.}}}
\newcommand{\cprob}{\xrightarrow{\Prob}}
\newcommand{\clp}[1]{\xrightarrow{L^{#1}}}
\newcommand{\cdist}{\xrightarrow{d}}

% --------------------------------------------------------------------------
% Custom commands -- common operators
% --------------------------------------------------------------------------
\DeclareMathOperator{\dom}{dom}
\DeclareMathOperator{\epi}{epi}
\DeclareMathOperator{\conv}{conv}
\DeclareMathOperator{\relint}{relint}
\DeclareMathOperator{\aff}{aff}
\DeclareMathOperator{\cl}{cl}
\DeclareMathOperator{\interior}{int}
\DeclareMathOperator{\tr}{tr}
\DeclareMathOperator{\diag}{diag}
\DeclareMathOperator{\rank}{rank}
\DeclareMathOperator{\sgn}{sgn}
\DeclareMathOperator{\supp}{supp}
\DeclareMathOperator{\argmin}{arg\,min}
\DeclareMathOperator{\argmax}{arg\,max}
\DeclareMathOperator{\prox}{prox}
\DeclareMathOperator{\dist}{dist}
\DeclareMathOperator{\diam}{diam}
\DeclareMathOperator{\Span}{span}

% --------------------------------------------------------------------------
% Custom commands -- linear algebra operators
% --------------------------------------------------------------------------
\DeclareMathOperator{\nullity}{nullity}
\DeclareMathOperator{\range}{range}
\DeclareMathOperator{\nullspace}{null}
\DeclareMathOperator{\proj}{proj}

% --------------------------------------------------------------------------
% Custom commands -- measure theory
% --------------------------------------------------------------------------
\newcommand{\powerset}{\mathcal{P}}
\newcommand{\Borel}{\mathcal{B}}

\pagestyle{fancy}
\fancyhf{}
\fancyhead[L]{\textsc{Lesson 6: Determinants and Eigenvalues}}
\fancyhead[R]{\thepage}
\title{Lesson 6: Determinants and Eigenvalues}
\author{AI/ML Mathematics}
\date{}
\begin{document}
\maketitle
\tableofcontents
\newpage

%% ============================================================
%% SECTION 1: MOTIVATION AND OVERVIEW
%% ============================================================
\section{Motivation and Overview}
\label{sec:motivation}

Consider a Markov chain with $n$ states and transition matrix
$P \in \R^{n \times n}$.  A fundamental question in reinforcement learning
is: how quickly does the chain converge to its stationary distribution?
The answer depends on the \emph{spectral gap} $1 - |\lambda_2|$, where
$\lambda_2$ is the second-largest eigenvalue of $P$ in absolute value.
But what exactly is an eigenvalue?  How do we know eigenvalues exist?
And how do we compute them?

The characteristic polynomial $\det(A - \lambda I) = 0$ is the bridge
connecting eigenvalues to the determinant.  This single equation governs
the convergence rate of iterative algorithms (value iteration, policy
iteration), the stability of gradient descent (all eigenvalues of the
Hessian must be positive for a local minimum), and the expressiveness of
linear transformations (invertibility is equivalent to a nonzero determinant).
The tools developed in this lesson will let us define determinants rigorously,
establish the eigenvalue--eigenvector framework, and prove the
diagonalizability criterion that underpins spectral methods throughout
machine learning.

\begin{rlconnection}[title={Where Determinants and Eigenvalues Appear in RL/ML}]
\begin{itemize}[nosep]
  \item \textbf{Spectral radius and Markov chain mixing.}
        The rate at which a Markov chain mixes to stationarity is governed by
        the spectral radius $\rho(P - \pi \mathbf{1}^\top)$; eigenvalue
        analysis is the primary tool for bounding this quantity.
  \item \textbf{Gradient descent convergence.}
        For quadratic objectives $f(x) = \tfrac{1}{2} x^\top A x - b^\top x$,
        the condition number $\kappa = \lambda_{\max}(A)/\lambda_{\min}(A)$
        determines the convergence rate of gradient descent.
  \item \textbf{Cayley--Hamilton and multi-step TD.}
        The Cayley--Hamilton theorem shows that $A^n$ can be expressed as a
        polynomial in $A$ of degree $< n$, which is used to analyze
        multi-step temporal difference methods and truncated matrix powers.
\end{itemize}
\end{rlconnection}

\paragraph{Prerequisites.}
Vector spaces, bases, dimension, linear maps, rank--nullity theorem,
matrix representation, change of basis, and similarity (Lesson~5).

\paragraph{Learning objectives.}
After completing this lesson, the student will be able to:
\begin{itemize}[nosep]
  \item \textbf{Define} the determinant via cofactor expansion and the
        Leibniz formula, and compute determinants of small matrices.
  \item \textbf{Prove} the multiplicativity of determinants
        ($\det(AB) = \det A \cdot \det B$) and the equivalence between
        invertibility and nonzero determinant.
  \item \textbf{Compute} eigenvalues and eigenvectors of matrices by
        solving the characteristic polynomial.
  \item \textbf{Derive} the relationship between algebraic and geometric
        multiplicity and apply the diagonalizability criterion.
  \item \textbf{State} and prove the Cayley--Hamilton theorem for
        matrices over $\F$.
\end{itemize}

%% ============================================================
%% SECTION 2: REFERENCES
%% ============================================================
\section{References}
\label{sec:references}

\begin{itemize}
  \item \citet[Chapter 10]{axler2015} -- Determinant theory developed after
        eigenvalues; provides a clean proof of the Cayley--Hamilton theorem.
  \item \citet[Chapter 5]{hoffman1971} -- Classical treatment of
        determinants via multilinear alternating forms.
  \item \citet[Chapter 5]{strang2016} -- Computational perspective on
        determinants with cofactor expansion and applications.
  \item \citet[Chapter 1]{horn2013} -- Eigenvalue theory, characteristic
        polynomial, and matrix factorizations.
  \item \citet[Chapter 6]{lax2007} -- Spectral theory with connections
        to operator theory and applications.
\end{itemize}

%% ============================================================
%% SECTION 3: DETERMINANTS
%% ============================================================
\section{Determinants}
\label{sec:determinants}

Throughout this lesson, $\F$ denotes either $\R$ or $\C$, and all matrices
are square unless stated otherwise.

\subsection{Permutations and the Sign Function}

\begin{definition}[Permutation and Sign]
\label{def:permutation}
A \emph{permutation} of $\{1, \ldots, n\}$ is a bijection
$\sigma\colon \{1, \ldots, n\} \to \{1, \ldots, n\}$.  The set of all
such permutations is denoted $S_n$.  A \emph{transposition} is a
permutation that swaps exactly two elements.  Every permutation can be
written as a composition of transpositions.  The \emph{sign} (or
\emph{signature}) of $\sigma$ is
\[
  \sgn(\sigma) = (-1)^{k},
\]
where $k$ is the number of transpositions in any decomposition of $\sigma$.
A permutation is \emph{even} if $\sgn(\sigma) = +1$ and \emph{odd} if
$\sgn(\sigma) = -1$.
\end{definition}

\begin{remark}
\label{rem:sign_well_defined}
The sign is well-defined: although a permutation can be written as a
product of transpositions in many ways, the parity of the number of
transpositions is invariant.  One proof uses the identity
$\sgn(\sigma) = \prod_{i < j} \frac{\sigma(j) - \sigma(i)}{j - i}$,
which is manifestly well-defined \citep[Chapter~5]{hoffman1971}.
\end{remark}

\begin{example}[Permutations of $\{1,2,3\}$]
\label{ex:S3}
The group $S_3$ has $3! = 6$ elements.  The identity $e = (1)(2)(3)$
has $\sgn(e) = +1$.  The transposition $\sigma = (1\;2)$ (swapping 1 and 2)
has $\sgn(\sigma) = -1$.  The cyclic permutation $\tau = (1\;2\;3)$
(mapping $1 \mapsto 2 \mapsto 3 \mapsto 1$) decomposes as
$(1\;3)(1\;2)$, so $\sgn(\tau) = +1$.
\end{example}

\subsection{The Leibniz Formula}

\begin{definition}[Determinant -- Leibniz Formula]
\label{def:det_leibniz}
Let $A = (a_{ij}) \in \F^{n \times n}$.  The \emph{determinant} of $A$ is
\[
  \det A = \sum_{\sigma \in S_n} \sgn(\sigma) \prod_{i=1}^{n} a_{i,\sigma(i)}.
\]
\end{definition}

\begin{example}[Determinants of $2 \times 2$ and $3 \times 3$ Matrices]
\label{ex:det_small}
For $A = \bigl(\begin{smallmatrix} a & b \\ c & d \end{smallmatrix}\bigr)$,
the two permutations of $\{1,2\}$ give
$\det A = ad - bc$.

For a $3 \times 3$ matrix $A = (a_{ij})$, expanding over all six
permutations in $S_3$ yields
\begin{align*}
  \det A &= a_{11}a_{22}a_{33} + a_{12}a_{23}a_{31} + a_{13}a_{21}a_{32} \\
         &\quad - a_{13}a_{22}a_{31} - a_{12}a_{21}a_{33} - a_{11}a_{23}a_{32}.
\end{align*}
\end{example}

\subsection{Cofactor Expansion}

\begin{definition}[Minor and Cofactor]
\label{def:minor_cofactor}
Let $A \in \F^{n \times n}$.  The \emph{$(i,j)$-minor} $M_{ij}$ is the
determinant of the $(n-1) \times (n-1)$ matrix obtained by deleting row $i$
and column $j$ of $A$.  The \emph{$(i,j)$-cofactor} is
$C_{ij} = (-1)^{i+j} M_{ij}$.
\end{definition}

\begin{theorem}[Cofactor Expansion]
\label{thm:cofactor_expansion}
For any fixed row $i$ or column $j$,
\[
  \det A = \sum_{j=1}^{n} a_{ij}\, C_{ij}
  \quad \text{(expansion along row $i$)},
\]
\[
  \det A = \sum_{i=1}^{n} a_{ij}\, C_{ij}
  \quad \text{(expansion along column $j$)}.
\]
\end{theorem}

\begin{proof}
We prove the row expansion.  In the Leibniz formula, fix row $i$ and
group terms by the column index $\sigma(i) = j$.  For each $j$, the
contribution from permutations with $\sigma(i) = j$ is
$a_{ij} \cdot (-1)^{i+j} \det A^{(ij)}$, where $A^{(ij)}$ is the matrix
with row $i$ and column $j$ deleted.  To verify the sign, note that
moving the factor $a_{i,\sigma(i)}$ from position $i$ in the product
to the front requires reindexing the remaining permutation on
$\{1,\ldots,n\} \setminus \{i\}$ to columns
$\{1,\ldots,n\} \setminus \{j\}$, which introduces a sign factor of
$(-1)^{i+j}$.  Summing over $j = 1, \ldots, n$ recovers
$\det A = \sum_j a_{ij} C_{ij}$.  The column expansion follows from
the transpose property (Proposition~\ref{prop:det_transpose}).
\end{proof}

\begin{example}[Cofactor Expansion of a $3 \times 3$ Matrix]
\label{ex:cofactor_3x3}
Let $A = \begin{pmatrix} 2 & 1 & 3 \\ 0 & -1 & 4 \\ 5 & 2 & 1 \end{pmatrix}$.
Expanding along row 1:
\begin{align*}
  \det A &= 2 \det\begin{pmatrix} -1 & 4 \\ 2 & 1 \end{pmatrix}
          - 1 \det\begin{pmatrix} 0 & 4 \\ 5 & 1 \end{pmatrix}
          + 3 \det\begin{pmatrix} 0 & -1 \\ 5 & 2 \end{pmatrix} \\
         &= 2(-1-8) - 1(0-20) + 3(0+5)
          = -18 + 20 + 15 = 17.
\end{align*}
\end{example}

\subsection{Properties of Determinants}

\begin{proposition}[Transpose]
\label{prop:det_transpose}
For any $A \in \F^{n \times n}$, $\det(A^\top) = \det A$.
\end{proposition}

\begin{proof}
In the Leibniz formula, $\det(A^\top) = \sum_{\sigma \in S_n} \sgn(\sigma)
\prod_{i=1}^n a_{\sigma(i),i}$.  Since $\sigma \mapsto \sigma^{-1}$ is a
bijection on $S_n$ with $\sgn(\sigma^{-1}) = \sgn(\sigma)$, reindexing
with $\tau = \sigma^{-1}$ gives
$\sum_{\tau \in S_n} \sgn(\tau) \prod_{i=1}^n a_{i,\tau(i)} = \det A$.
\end{proof}

\begin{proposition}[Determinant and Row Operations]
\label{prop:det_row_ops}
Let $A \in \F^{n \times n}$.
\begin{enumerate}[nosep]
  \item \textbf{Multilinearity:} The determinant is linear in each row
        (holding the other rows fixed).  That is, if row $i$ of $A$ is
        $\alpha r + \beta r'$ and $B, C$ are obtained by replacing row $i$
        with $r$ and $r'$ respectively, then $\det A = \alpha \det B + \beta \det C$.
  \item \textbf{Alternating:} If two rows of $A$ are interchanged, the
        determinant changes sign.  In particular, if two rows are equal,
        $\det A = 0$.
  \item \textbf{Row addition:} Adding a scalar multiple of one row to
        another does not change the determinant.
  \item \textbf{Row scaling:} Multiplying a single row by $\alpha \in \F$
        multiplies the determinant by $\alpha$.
\end{enumerate}
\end{proposition}

\begin{proof}
(1) follows directly from the Leibniz formula: each term
$\sgn(\sigma) \prod_i a_{i,\sigma(i)}$ is linear in the entries of any
fixed row.

(2) Let $A'$ be obtained from $A$ by swapping rows $i$ and $j$.  This is
equivalent to composing each permutation $\sigma$ with the transposition
$(i\;j)$, which changes the sign.  If rows $i$ and $j$ are equal, then
$A = A'$ and $\det A = -\det A$, so $\det A = 0$.

(3) Adding $\alpha$ times row $j$ to row $i$ gives, by multilinearity,
$\det A' = \det A + \alpha \det B$, where $B$ has row $j$ appearing in
both positions $i$ and $j$.  By (2), $\det B = 0$.

(4) is immediate from (1) with $r' = 0$.
\end{proof}

\begin{example}[Determinant via Row Reduction]
\label{ex:det_row_reduction}
Let $A = \begin{pmatrix} 1 & 2 & 3 \\ 2 & 5 & 7 \\ 3 & 7 & 10 \end{pmatrix}$.
Subtract 2 times row 1 from row 2, and 3 times row 1 from row 3:
$\begin{pmatrix} 1 & 2 & 3 \\ 0 & 1 & 1 \\ 0 & 1 & 1 \end{pmatrix}$.
Two rows are now equal, so $\det A = 0$.  By the criterion we will prove
shortly, $A$ is not invertible.
\end{example}

\subsection{Multiplicativity of the Determinant}

\begin{theorem}[Multiplicativity]
\label{thm:det_mult}
For $A, B \in \F^{n \times n}$,
\[
  \det(AB) = \det A \cdot \det B.
\]
\end{theorem}

\begin{proof}
\textbf{Case 1:} $A$ is not invertible.  Then $\rank(A) < n$, so
$\rank(AB) \leq \rank(A) < n$.  Hence $AB$ is not invertible, and we
will show in Theorem~\ref{thm:det_invertibility} that
$\det(AB) = 0 = \det A \cdot \det B$.

\textbf{Case 2:} $A$ is invertible.  Every invertible matrix is a product
of elementary matrices (each representing a single row operation).
An elementary matrix $E$ satisfies $\det(EB) = (\det E)(\det B)$ by
Proposition~\ref{prop:det_row_ops}: swapping two rows multiplies the
determinant by $-1 = \det E$; adding a multiple of one row to another
multiplies by $1 = \det E$; scaling a row by $\alpha$ multiplies by
$\alpha = \det E$.  Writing $A = E_1 \cdots E_k$, repeated application
gives
\[
  \det(AB) = \det(E_1 \cdots E_k B)
  = (\det E_1) \cdots (\det E_k)(\det B)
  = \det(E_1 \cdots E_k) \det B
  = \det A \cdot \det B. \qedhere
\]
\end{proof}

\subsection{Determinants and Invertibility}

\begin{keyresult}[title={Determinant--Invertibility Equivalence}]
\begin{theorem}[Determinant and Invertibility]
\label{thm:det_invertibility}
A matrix $A \in \F^{n \times n}$ is invertible if and only if $\det A \neq 0$.
When $A$ is invertible, $\det(A^{-1}) = (\det A)^{-1}$.
\end{theorem}
\end{keyresult}

\begin{proof}
$(\Rightarrow)$\; If $A$ is invertible, then $AA^{-1} = I$, so by
multiplicativity $\det A \cdot \det(A^{-1}) = \det I = 1$.  In particular,
$\det A \neq 0$ and $\det(A^{-1}) = (\det A)^{-1}$.

$(\Leftarrow)$\; If $A$ is not invertible, then $\rank(A) < n$, so the
rows of $A$ are linearly dependent.  Row reduction (which preserves
$\det = 0$ or $\det \neq 0$) produces at least one zero row.  A matrix
with a zero row has determinant $0$ (every term in the Leibniz formula
contains a factor from this row, which is $0$).  Hence $\det A = 0$.
\end{proof}

\begin{example}[Invertibility Check]
\label{ex:invertibility}
For $A = \begin{pmatrix} 1 & 2 \\ 3 & 6 \end{pmatrix}$, we have
$\det A = 6 - 6 = 0$, so $A$ is singular.  Indeed, row 2 is 3 times row 1.
For $B = \begin{pmatrix} 1 & 2 \\ 3 & 7 \end{pmatrix}$, we have
$\det B = 7 - 6 = 1 \neq 0$, so $B$ is invertible.
\end{example}

The determinant provides a scalar invariant that detects invertibility.
We now turn to the richer structure revealed by the eigenvalue problem.

%% ============================================================
%% SECTION 4: EIGENVALUES AND EIGENVECTORS
%% ============================================================
\section{Eigenvalues and Eigenvectors}
\label{sec:eigenvalues}

\subsection{Basic Definitions}

\begin{definition}[Eigenvalue and Eigenvector]
\label{def:eigenvalue}
Let $A \in \F^{n \times n}$.  A scalar $\lambda \in \F$ is an
\emph{eigenvalue} of $A$ if there exists a nonzero vector $v \in \F^n$
such that $Av = \lambda v$.  Such a vector $v$ is called an
\emph{eigenvector} of $A$ corresponding to $\lambda$.
\end{definition}

\begin{definition}[Eigenspace]
\label{def:eigenspace}
The \emph{eigenspace} of $A$ corresponding to eigenvalue $\lambda$ is
\[
  E_\lambda = \nullspace(A - \lambda I) = \{v \in \F^n : Av = \lambda v\}.
\]
This is a subspace of $\F^n$ (it is the null space of the linear map
$A - \lambda I$).  The eigenspace includes the zero vector, but by
convention eigenvectors are nonzero.
\end{definition}

\begin{example}[Eigenvalues of a $2 \times 2$ Matrix]
\label{ex:eigen_2x2}
Let $A = \begin{pmatrix} 4 & 1 \\ 2 & 3 \end{pmatrix}$.  If $Av = \lambda v$
for nonzero $v$, then $(A - \lambda I)v = 0$ has a nontrivial solution,
which requires $\det(A - \lambda I) = 0$.  We compute
\[
  \det(A - \lambda I)
  = \det\begin{pmatrix} 4-\lambda & 1 \\ 2 & 3-\lambda \end{pmatrix}
  = (4-\lambda)(3-\lambda) - 2 = \lambda^2 - 7\lambda + 10
  = (\lambda - 2)(\lambda - 5).
\]
The eigenvalues are $\lambda_1 = 2$ and $\lambda_2 = 5$.  For $\lambda_1 = 2$:
$(A - 2I)v = 0$ gives $2v_1 + v_2 = 0$, so $E_2 = \Span\{(1,-2)\}$.
For $\lambda_2 = 5$: $-v_1 + v_2 = 0$, so $E_5 = \Span\{(1,1)\}$.
\end{example}

\subsection{The Characteristic Polynomial}

\begin{definition}[Characteristic Polynomial]
\label{def:char_poly}
The \emph{characteristic polynomial} of $A \in \F^{n \times n}$ is
\[
  p_A(\lambda) = \det(A - \lambda I).
\]
This is a polynomial of degree $n$ in $\lambda$ with leading coefficient
$(-1)^n$.  The eigenvalues of $A$ are precisely the roots of $p_A$.
\end{definition}

\begin{remark}
\label{rem:char_poly_similarity}
Similar matrices have the same characteristic polynomial: if $B = P^{-1}AP$,
then $\det(B - \lambda I) = \det(P^{-1}(A - \lambda I)P)
= \det(P^{-1}) \det(A - \lambda I) \det(P) = \det(A - \lambda I)$.
Hence eigenvalues are invariants of the linear operator, not of a
particular matrix representation.
\end{remark}

\begin{proposition}[Trace and Determinant from Eigenvalues]
\label{prop:trace_det_eigen}
If $A \in \F^{n \times n}$ has eigenvalues $\lambda_1, \ldots, \lambda_n$
(counted with algebraic multiplicity, over the algebraic closure of $\F$),
then
\[
  \tr(A) = \sum_{i=1}^{n} \lambda_i, \qquad
  \det(A) = \prod_{i=1}^{n} \lambda_i.
\]
\end{proposition}

\begin{proof}
The characteristic polynomial factors (over the algebraic closure) as
$p_A(\lambda) = (-1)^n(\lambda - \lambda_1) \cdots (\lambda - \lambda_n)$.
Setting $\lambda = 0$: $\det A = p_A(0) = (-1)^n(-\lambda_1) \cdots (-\lambda_n)
= \prod_i \lambda_i$.
Comparing the coefficient of $\lambda^{n-1}$: in $\det(A - \lambda I)$,
the $\lambda^{n-1}$ term arises only from the diagonal product
$(a_{11} - \lambda) \cdots (a_{nn} - \lambda)$, giving coefficient
$(-1)^{n-1}(a_{11} + \cdots + a_{nn}) = (-1)^{n-1} \tr(A)$.
From the factored form, the $\lambda^{n-1}$ coefficient is
$(-1)^{n-1}(\lambda_1 + \cdots + \lambda_n)$.  Equating gives
$\tr(A) = \sum_i \lambda_i$.
\end{proof}

\begin{rlconnection}[title={Spectral Radius and Markov Chain Mixing}]
For a stochastic matrix $P$ (rows sum to 1, nonneg entries), $\lambda_1 = 1$
is always an eigenvalue.  The \emph{spectral radius} of $P$ restricted
to the complement of the stationary distribution is
$\rho = \max_{i \geq 2} |\lambda_i|$.  The mixing time satisfies
$t_{\mathrm{mix}} = \Theta(1/(1-\rho))$.  In policy evaluation for RL,
faster mixing means that Monte Carlo rollouts more quickly yield
representative samples, reducing variance.
\end{rlconnection}

\subsection{Algebraic and Geometric Multiplicity}

\begin{definition}[Algebraic and Geometric Multiplicity]
\label{def:multiplicity}
Let $\lambda$ be an eigenvalue of $A$.
\begin{itemize}[nosep]
  \item The \emph{algebraic multiplicity} $\mathrm{am}(\lambda)$ is the
        multiplicity of $\lambda$ as a root of the characteristic polynomial
        $p_A$.
  \item The \emph{geometric multiplicity} $\mathrm{gm}(\lambda)$ is the
        dimension of the eigenspace $E_\lambda = \nullspace(A - \lambda I)$.
\end{itemize}
\end{definition}

\begin{theorem}[Multiplicity Inequality]
\label{thm:multiplicity_ineq}
For every eigenvalue $\lambda$ of $A \in \F^{n \times n}$,
\[
  1 \leq \mathrm{gm}(\lambda) \leq \mathrm{am}(\lambda).
\]
\end{theorem}

\begin{proof}
The lower bound $\mathrm{gm}(\lambda) \geq 1$ holds because the
eigenspace contains at least one eigenvector (by definition of eigenvalue).

For the upper bound, let $k = \mathrm{gm}(\lambda)$ and let
$\{v_1, \ldots, v_k\}$ be a basis of $E_\lambda$.  Extend this to a basis
$\{v_1, \ldots, v_k, v_{k+1}, \ldots, v_n\}$ of $\F^n$.  Let $P$ be the
invertible matrix with columns $v_1, \ldots, v_n$.  Then
\[
  P^{-1}AP = \begin{pmatrix} \lambda I_k & * \\ 0 & C \end{pmatrix},
\]
because $Av_i = \lambda v_i$ for $i = 1, \ldots, k$.  The characteristic
polynomial is
\[
  p_A(t) = \det(P^{-1}AP - tI)
  = (\lambda - t)^k \det(C - tI).
\]
Hence $(\lambda - t)^k$ divides $p_A(t)$, so
$\mathrm{am}(\lambda) \geq k = \mathrm{gm}(\lambda)$.
\end{proof}

\begin{example}[Strict Inequality of Multiplicities]
\label{ex:strict_mult}
Consider $A = \begin{pmatrix} 2 & 1 \\ 0 & 2 \end{pmatrix}$.  The
characteristic polynomial is $(2 - \lambda)^2$, so $\lambda = 2$ has
$\mathrm{am}(2) = 2$.  But $A - 2I = \begin{pmatrix} 0 & 1 \\ 0 & 0 \end{pmatrix}$,
which has $\nullity = 1$, so $\mathrm{gm}(2) = 1 < 2 = \mathrm{am}(2)$.
\end{example}

\subsection{Independence of Eigenvectors for Distinct Eigenvalues}

\begin{theorem}[Eigenvectors for Distinct Eigenvalues Are Independent]
\label{thm:distinct_eigen_indep}
Let $\lambda_1, \ldots, \lambda_k$ be distinct eigenvalues of
$A \in \F^{n \times n}$ with corresponding eigenvectors $v_1, \ldots, v_k$.
Then $v_1, \ldots, v_k$ are linearly independent.
\end{theorem}

\begin{proof}
We proceed by induction on $k$.  The base case $k = 1$ is immediate:
a single eigenvector is nonzero, hence independent.

For the inductive step, suppose the result holds for $k - 1$ distinct
eigenvalues.  Assume
\begin{equation}
\label{eq:eigen_dep}
  c_1 v_1 + c_2 v_2 + \cdots + c_k v_k = 0.
\end{equation}
Apply $A$ to both sides:
\begin{equation}
\label{eq:eigen_dep_A}
  c_1 \lambda_1 v_1 + c_2 \lambda_2 v_2 + \cdots + c_k \lambda_k v_k = 0.
\end{equation}
Multiply \eqref{eq:eigen_dep} by $\lambda_k$ and subtract from
\eqref{eq:eigen_dep_A}:
\[
  c_1(\lambda_1 - \lambda_k) v_1 + c_2(\lambda_2 - \lambda_k) v_2
  + \cdots + c_{k-1}(\lambda_{k-1} - \lambda_k) v_{k-1} = 0.
\]
By the inductive hypothesis, $v_1, \ldots, v_{k-1}$ are independent.
Since all $\lambda_i - \lambda_k \neq 0$ for $i < k$, we conclude
$c_i = 0$ for $i = 1, \ldots, k-1$.  Substituting back into
\eqref{eq:eigen_dep} gives $c_k v_k = 0$, and since $v_k \neq 0$,
$c_k = 0$.
\end{proof}

\begin{rlconnection}[title={Eigenvalues and Gradient Descent Convergence}]
For a quadratic objective $f(x) = \frac{1}{2}x^\top A x - b^\top x$ with
$A$ symmetric positive definite, gradient descent with step size $\alpha$
has update $x_{k+1} = (I - \alpha A)x_k + \alpha b$.  The error
$e_k = x_k - x^*$ satisfies $e_k = (I - \alpha A)^k e_0$.  Since the
eigenvectors of $A$ (which are real and independent by
Theorem~\ref{thm:distinct_eigen_indep} when eigenvalues are distinct)
form a basis, each component of $e_k$ in this eigenbasis decays as
$(1 - \alpha \lambda_i)^k$.  Convergence requires
$|1 - \alpha \lambda_i| < 1$ for all $i$, optimized by
$\alpha^* = 2/(\lambda_{\min} + \lambda_{\max})$ with rate
$(\kappa - 1)/(\kappa + 1)$ where $\kappa = \lambda_{\max}/\lambda_{\min}$.
\end{rlconnection}

Having established when eigenvalues exist and that eigenvectors for
distinct eigenvalues are independent, we now ask: when can we find a
basis of eigenvectors for the entire space?

%% ============================================================
%% SECTION 5: DIAGONALIZABILITY
%% ============================================================
\section{Diagonalizability}
\label{sec:diagonalizability}

\begin{definition}[Diagonalizable Matrix]
\label{def:diagonalizable}
A matrix $A \in \F^{n \times n}$ is \emph{diagonalizable} if there exists
an invertible matrix $P \in \F^{n \times n}$ such that $P^{-1}AP$ is
diagonal.  Equivalently, $A$ is diagonalizable if and only if $\F^n$ has a
basis consisting of eigenvectors of $A$.
\end{definition}

\begin{keyresult}[title={Diagonalizability Criterion}]
\begin{theorem}[Diagonalizability Criterion]
\label{thm:diag_criterion}
A matrix $A \in \F^{n \times n}$ is diagonalizable if and only if both
of the following conditions hold:
\begin{enumerate}[nosep]
  \item The characteristic polynomial $p_A$ splits completely over $\F$
        (i.e., all eigenvalues lie in $\F$).
  \item For every eigenvalue $\lambda$,
        $\mathrm{gm}(\lambda) = \mathrm{am}(\lambda)$.
\end{enumerate}
\end{theorem}
\end{keyresult}

\begin{proof}
$(\Rightarrow)$\; Suppose $A$ is diagonalizable: $P^{-1}AP = D$ where
$D = \diag(\lambda_1, \ldots, \lambda_n)$ with $\lambda_i \in \F$.  Then
$p_A(\lambda) = p_D(\lambda) = \prod_{i=1}^n (\lambda_i - \lambda)$,
which splits over $\F$.  If $\mu$ appears $m$ times on the diagonal of $D$,
then $\mathrm{am}(\mu) = m$.  The columns of $P$ corresponding to diagonal
entries equal to $\mu$ are $m$ linearly independent eigenvectors for $\mu$,
so $\mathrm{gm}(\mu) \geq m = \mathrm{am}(\mu)$.  Combined with
Theorem~\ref{thm:multiplicity_ineq}, $\mathrm{gm}(\mu) = \mathrm{am}(\mu)$.

$(\Leftarrow)$\; Let $\lambda_1, \ldots, \lambda_r$ be the distinct
eigenvalues of $A$ with algebraic multiplicities $m_1, \ldots, m_r$
summing to $n$.  By assumption, each eigenspace $E_{\lambda_i}$ has
dimension $m_i$.  We claim the union of bases of all eigenspaces forms a
basis of $\F^n$.

The total count is $m_1 + \cdots + m_r = n$, so it suffices to show
linear independence.  Suppose
$\sum_{i=1}^r w_i = 0$ where $w_i \in E_{\lambda_i}$.  Write each $w_i$
as a combination of a basis of $E_{\lambda_i}$.  The vectors from
different eigenspaces lie in distinct eigenspaces; by an extension of the
argument in Theorem~\ref{thm:distinct_eigen_indep} (eigenvectors from
distinct eigenspaces are independent), each $w_i = 0$.  Since the basis
vectors within each $E_{\lambda_i}$ are independent, all coefficients
vanish.  Hence these $n$ eigenvectors form a basis of $\F^n$, and $A$ is
diagonalizable.
\end{proof}

\begin{corollary}
\label{cor:distinct_diag}
If $A \in \F^{n \times n}$ has $n$ distinct eigenvalues in $\F$, then
$A$ is diagonalizable.
\end{corollary}

\begin{proof}
Each eigenvalue has algebraic multiplicity 1, and geometric multiplicity
is at least 1 by Theorem~\ref{thm:multiplicity_ineq}.  Hence
$\mathrm{gm}(\lambda_i) = 1 = \mathrm{am}(\lambda_i)$ for each $i$,
and the criterion is satisfied.
\end{proof}

\begin{warning}[title={Not Every Matrix Is Diagonalizable}]
The matrix $A = \begin{pmatrix} 0 & 1 \\ 0 & 0 \end{pmatrix}$ has
characteristic polynomial $\lambda^2$, so $\lambda = 0$ is the only
eigenvalue with $\mathrm{am}(0) = 2$.  But $\rank(A - 0 \cdot I) = 1$,
so $\mathrm{gm}(0) = 1 < 2$.  Hence $A$ is not diagonalizable.

Over $\R$, even $2 \times 2$ matrices may fail: the rotation matrix
$R = \begin{pmatrix} 0 & -1 \\ 1 & 0 \end{pmatrix}$ has characteristic
polynomial $\lambda^2 + 1$, which has no real roots.  Over $\C$, the
eigenvalues are $\pm i$ (distinct), so $R$ is diagonalizable over $\C$
but not over $\R$.
\end{warning}

\begin{example}[Diagonalization Procedure]
\label{ex:diagonalization}
Let $A = \begin{pmatrix} 1 & 2 \\ 0 & 3 \end{pmatrix}$.  The
characteristic polynomial is $(1 - \lambda)(3 - \lambda)$, giving
distinct eigenvalues $\lambda_1 = 1$ and $\lambda_2 = 3$.

For $\lambda_1 = 1$: $(A - I)v = 0$ gives
$\begin{pmatrix} 0 & 2 \\ 0 & 2 \end{pmatrix}v = 0$, so
$v_1 = (1, 0)^\top$.

For $\lambda_2 = 3$: $(A - 3I)v = 0$ gives
$\begin{pmatrix} -2 & 2 \\ 0 & 0 \end{pmatrix}v = 0$, so
$v_2 = (1, 1)^\top$.

Setting $P = \begin{pmatrix} 1 & 1 \\ 0 & 1 \end{pmatrix}$, we get
$P^{-1}AP = \begin{pmatrix} 1 & 0 \\ 0 & 3 \end{pmatrix}$.
\end{example}

With diagonalizability understood, we conclude with the Cayley--Hamilton
theorem, which provides a universal polynomial identity satisfied by every
square matrix.

%% ============================================================
%% SECTION 6: THE CAYLEY-HAMILTON THEOREM
%% ============================================================
\section{The Cayley--Hamilton Theorem}
\label{sec:cayley_hamilton}

\begin{theorem}[Cayley--Hamilton]
\label{thm:cayley_hamilton}
Let $A \in \F^{n \times n}$ and let $p_A(\lambda) = \det(A - \lambda I)$
be its characteristic polynomial.  Then $p_A(A) = 0$ (the zero matrix).
That is, if $p_A(\lambda) = c_0 + c_1 \lambda + \cdots + c_{n-1}\lambda^{n-1}
+ (-1)^n \lambda^n$, then
\[
  c_0 I + c_1 A + \cdots + c_{n-1} A^{n-1} + (-1)^n A^n = 0.
\]
\end{theorem}

\begin{proof}
Let $\mathrm{adj}(A - \lambda I) = B_0 + B_1 \lambda + \cdots + B_{n-1} \lambda^{n-1}$
be the adjugate, whose entries are cofactors (polynomials of degree
$\leq n-1$ in $\lambda$).  The cofactor-adjugate identity gives
\[
  (A - \lambda I)(B_0 + B_1 \lambda + \cdots + B_{n-1}\lambda^{n-1})
  = p_A(\lambda) \cdot I.
\]
Write $p_A(\lambda) = c_0 + c_1\lambda + \cdots + (-1)^n \lambda^n$.
Equating coefficients of $\lambda^k$ for $k = 0, \ldots, n$ yields
$n+1$ matrix equations.  Multiplying the $k$-th equation on the left
by $A^k$ and summing, the left side telescopes (each $A^k B_{k-1}$
appears once with $+1$ and once with $-1$), giving
\[
  0 = c_0 I + c_1 A + c_2 A^2 + \cdots + (-1)^n A^n = p_A(A). \qedhere
\]
\end{proof}

\begin{example}[Cayley--Hamilton for $2 \times 2$]
\label{ex:cayley_hamilton_2x2}
Let $A = \begin{pmatrix} 1 & 3 \\ 2 & 4 \end{pmatrix}$.  The characteristic
polynomial is $p_A(\lambda) = \lambda^2 - 5\lambda - 2$.  The theorem
asserts $A^2 - 5A - 2I = 0$.  We verify:
\[
  A^2 = \begin{pmatrix} 7 & 15 \\ 10 & 22 \end{pmatrix}, \quad
  5A = \begin{pmatrix} 5 & 15 \\ 10 & 20 \end{pmatrix}, \quad
  A^2 - 5A - 2I = \begin{pmatrix} 0 & 0 \\ 0 & 0 \end{pmatrix}. \quad \checkmark
\]
\end{example}

\begin{example}[Expressing $A^{-1}$ via Cayley--Hamilton]
\label{ex:ch_inverse}
If $A$ is invertible and $p_A(\lambda) = c_0 + c_1\lambda + \cdots + (-1)^n\lambda^n$
with $c_0 = \det A \neq 0$, then $p_A(A) = 0$ gives
\[
  c_0 I + c_1 A + \cdots + (-1)^n A^n = 0.
\]
Multiplying by $A^{-1}$ and solving:
\[
  A^{-1} = -\frac{1}{c_0}\bigl(c_1 I + c_2 A + \cdots + (-1)^n A^{n-1}\bigr).
\]
This expresses $A^{-1}$ as a polynomial in $A$ of degree $n-1$.
\end{example}

\begin{rlconnection}[title={Cayley--Hamilton in Multi-Step TD}]
In temporal difference learning with eligibility traces, the
$\lambda$-return involves the matrix geometric series
$\sum_{k=0}^{\infty} (\gamma P^\pi)^k$ where $P^\pi$ is the transition
matrix.  By the Cayley--Hamilton theorem, $A^n$ (and all higher powers)
can be expressed as linear combinations of $I, A, \ldots, A^{n-1}$.
This means the effective ``memory'' of the matrix powers is at most $n$
steps, and in practice, truncated sums at horizon $n$ capture the full
behavior.  This observation justifies efficient implementations of
multi-step TD methods in finite-state MDPs.
\end{rlconnection}

%% ============================================================
%% SECTION 7: EXERCISES
%% ============================================================
\section{Exercises}
\label{sec:exercises}

\begin{exercise}[Determinant Computation]
\label{exer:det_computation}
\textnormal{(\(*\))} \\
Compute the determinant of each matrix:
\begin{enumerate}[nosep,label=(\alph*)]
  \item $A = \begin{pmatrix} 3 & 1 & -2 \\ 0 & 4 & 5 \\ 1 & 0 & 3 \end{pmatrix}$.
  \item $B = \begin{pmatrix} 1 & 2 & 3 & 4 \\ 0 & 1 & 2 & 3 \\
        0 & 0 & 1 & 2 \\ 0 & 0 & 0 & 1 \end{pmatrix}$.
\end{enumerate}
\end{exercise}

\begin{exercise}[Eigenvalue Analysis]
\label{exer:eigenvalue_analysis}
\textnormal{(\(*\))} \\
Let $A = \begin{pmatrix} 2 & 1 & 0 \\ 0 & 2 & 0 \\ 0 & 0 & 3 \end{pmatrix}$.
Find all eigenvalues, their algebraic and geometric multiplicities, and
determine whether $A$ is diagonalizable.
\end{exercise}

\begin{exercise}[Determinant Identities]
\label{exer:det_identities}
\textnormal{(\(**\))} \\
Let $A, B \in \F^{n \times n}$ with $A$ invertible.
\begin{enumerate}[nosep,label=(\alph*)]
  \item Prove that $\det(A^{-1}BA) = \det B$.
  \item Prove that $A$ and $A^\top$ have the same eigenvalues.
  \item Show by example that $A$ and $A^\top$ can have different
        eigenvectors.
\end{enumerate}
\end{exercise}

\begin{exercise}[Markov Chain Spectral Analysis]
\label{exer:markov_spectral}
\textnormal{(\(**\))} \\
A Markov chain on states $\{1, 2, 3\}$ has transition matrix
$P = \begin{pmatrix} 0.5 & 0.3 & 0.2 \\ 0.2 & 0.5 & 0.3 \\
0.3 & 0.2 & 0.5 \end{pmatrix}$.
\begin{enumerate}[nosep,label=(\alph*)]
  \item Verify that $\lambda_1 = 1$ is an eigenvalue and find a
        corresponding left eigenvector (the stationary distribution).
  \item Compute all eigenvalues of $P$ and determine the spectral gap
        $1 - |\lambda_2|$.
  \item Use the spectral gap to bound the number of steps needed for the
        distribution to be within $\epsilon = 0.01$ of stationarity.
\end{enumerate}
\textit{Hint:} Row sums equal 1, so $(1,1,1)^\top$ is always an
eigenvector for $\lambda = 1$.
\end{exercise}

\begin{exercise}[Cayley--Hamilton Application]
\label{exer:cayley_hamilton}
\textnormal{(\(**\))} \\
Let $A = \begin{pmatrix} 1 & 1 \\ 0 & 2 \end{pmatrix}$.
\begin{enumerate}[nosep,label=(\alph*)]
  \item Compute the characteristic polynomial $p_A(\lambda)$.
  \item Verify $p_A(A) = 0$ by direct computation.
  \item Use the Cayley--Hamilton theorem to express $A^{-1}$ as
        $\alpha I + \beta A$ for explicit $\alpha, \beta$.
  \item Use the result to compute $A^5$ as $\gamma I + \delta A$
        for explicit $\gamma, \delta$.
\end{enumerate}
\end{exercise}

\begin{exercise}[Diagonalizability Proof]
\label{exer:diag_proof}
\textnormal{(\(***\))} \\
Let $A, B \in \F^{n \times n}$ be diagonalizable.  Prove or disprove:
\begin{enumerate}[nosep,label=(\alph*)]
  \item $A + B$ is diagonalizable.
  \item $AB$ is diagonalizable.
  \item If $A$ and $B$ are \emph{simultaneously} diagonalizable (i.e.,
        there exists a single invertible $P$ such that both $P^{-1}AP$ and
        $P^{-1}BP$ are diagonal), then $AB = BA$.  Prove this and state the
        converse.
\end{enumerate}
\textit{Hint:} For (a) and (b), consider $2 \times 2$ counterexamples.
For (c), use the fact that diagonal matrices commute.
\end{exercise}

%% ============================================================
%% SECTION 8: SUMMARY
%% ============================================================
\section{Summary}
\label{sec:summary}

\begin{itemize}
  \item The \emph{determinant} $\det A$ is defined by the Leibniz formula
        $\sum_{\sigma \in S_n} \sgn(\sigma) \prod_i a_{i,\sigma(i)}$ and
        can be computed recursively via cofactor expansion along any row or
        column.

  \item The determinant is \emph{multilinear} and \emph{alternating} in
        the rows; it is invariant under row addition and changes sign under
        row swaps.

  \item $\det(AB) = \det A \cdot \det B$ (multiplicativity), and $A$ is
        invertible if and only if $\det A \neq 0$.

  \item An \emph{eigenvalue} $\lambda$ of $A$ satisfies $Av = \lambda v$
        for some nonzero $v$; equivalently, $\lambda$ is a root of the
        \emph{characteristic polynomial} $p_A(\lambda) = \det(A - \lambda I)$.

  \item The \emph{trace} equals the sum of eigenvalues and the
        \emph{determinant} equals the product of eigenvalues.

  \item For each eigenvalue, $1 \leq \mathrm{gm}(\lambda) \leq
        \mathrm{am}(\lambda)$.  Eigenvectors for distinct eigenvalues
        are linearly independent.

  \item $A$ is \emph{diagonalizable} if and only if $p_A$ splits over
        $\F$ and $\mathrm{gm}(\lambda) = \mathrm{am}(\lambda)$ for every
        eigenvalue $\lambda$.  A sufficient condition is that $A$ has $n$
        distinct eigenvalues.

  \item The \emph{Cayley--Hamilton theorem} states $p_A(A) = 0$: every
        matrix satisfies its own characteristic polynomial.  This allows
        expressing high powers $A^k$ (for $k \geq n$) as polynomials
        in $A$ of degree at most $n - 1$.
\end{itemize}

%% ============================================================
%% SECTION 9: FURTHER READING
%% ============================================================
\section{Further Reading}
\label{sec:further_reading}

\begin{itemize}
  \item \textbf{Lesson 7} (Inner Products and Spectral Theory) introduces
        inner product structure and proves the spectral theorem for real
        symmetric matrices: every such matrix is orthogonally
        diagonalizable.  This strengthens the diagonalizability criterion
        by providing an orthonormal eigenbasis.

  \item \textbf{Lesson 8} (SVD and Applications) combines eigenvalue theory
        with inner products to derive the singular value decomposition,
        the most important matrix factorization in data science and ML.

  \item The Jordan normal form (not covered here) provides a canonical
        form for non-diagonalizable matrices.  See \citet[Chapter~8]{hoffman1971}
        or \citet[Chapter~3]{horn2013}.

  \item \citet[Chapter~1]{horn2013} gives a thorough treatment of
        eigenvalue inequalities (Weyl, Cauchy interlacing) used in
        matrix perturbation theory, which underlies the analysis of
        approximate methods in RL.

  \item For the historical development of determinants and eigenvalue
        theory, see \citet{lax2007}, which traces the ideas from Cayley
        and Sylvester through modern operator theory.
\end{itemize}

\bibliography{linear-algebra-refs}
\end{document}
